\documentclass[11pt,a4paper]{article}

\usepackage[utf8]{inputenc}
\usepackage{amsmath,amssymb,amsthm}
\usepackage{hyperref}
\usepackage{geometry}
\geometry{margin=1in}

\newcommand{\R}{\mathbb{R}}
\newcommand{\Z}{\mathbb{Z}}
\newcommand{\C}{\mathbb{C}}
\newtheorem{conjecture}{Conjecture}
\newtheorem{proposition}{Proposition}
\newtheorem{remark}{Remark}

\title{The Poisson Algebra of the Dyson Log-Gas:\\
  A Bridge Between Gravitational Algebras and Prime Distribution}
\author{B.\ Shepp}
\date{April 2026}

\begin{document}
\maketitle

\begin{abstract}
We establish a direct connection between the Poisson bracket algebra
generated by pairwise Hamiltonians in the $N$-body problem and the
statistical mechanics of Riemann zeta zeros. The Montgomery--Odlyzko law
identifies zeta zero correlations with GUE eigenvalue statistics, whose
joint density is the Boltzmann weight of the Dyson logarithmic Coulomb gas.
This gas decomposes into pairwise Hamiltonians $H_{ij} = -\log|t_i - t_j|$
plus harmonic confinement---precisely the system computed by the
\texttt{NBodyAlgebra} engine. The universality of the dimension sequence
$[3, 6, 17, 116]$ across singular potentials implies that the algebraic
structure governing zeta zero correlations belongs to the same universality
class as Newtonian gravity.
\end{abstract}

\tableofcontents

%% ====================================================================
\section{The Montgomery--Odlyzko Law}
\label{sec:montgomery}

Let $\rho_n = \tfrac{1}{2} + i\gamma_n$ denote the nontrivial zeros of
the Riemann zeta function $\zeta(s)$, ordered by imaginary part
$0 < \gamma_1 \le \gamma_2 \le \cdots$. Montgomery~\cite{Montgomery1973}
conjectured, and Odlyzko~\cite{Odlyzko1987} numerically verified to high
precision, that the pair correlation function of the normalized spacings
$\tilde\gamma_n = \gamma_n \cdot \frac{\log(\gamma_n / 2\pi e)}{2\pi}$
converges to
\begin{equation}\label{eq:pair-correlation}
  R_2(\alpha) = 1 - \left(\frac{\sin \pi\alpha}{\pi\alpha}\right)^2,
\end{equation}
which is exactly the pair correlation of eigenvalues of random matrices
drawn from the Gaussian Unitary Ensemble (GUE).

This connection extends beyond pair statistics. The $k$-point correlation
functions of zeta zeros match those of GUE for all $k$ tested numerically
\cite{Rudnick1996, Hejhal1994}. We now make this correspondence
Hamiltonian.

%% ====================================================================
\section{GUE as a Dyson Log-Gas}
\label{sec:gue}

The GUE is the probability measure on $N \times N$ Hermitian matrices
$M$ proportional to $\exp(-\tfrac{N}{2}\operatorname{Tr}M^2)$. Its
eigenvalue joint density is
\begin{equation}\label{eq:gue-density}
  P(t_1, \ldots, t_N) = C_N \prod_{i<j} |t_i - t_j|^2
  \cdot \exp\!\left(-\frac{1}{2}\sum_{i=1}^N t_i^2\right),
\end{equation}
where $C_N$ is a normalization constant (the Selberg integral).

\subsection{Boltzmann interpretation}

Writing $P = C_N \, e^{-\beta E}$ with $\beta = 1$, we identify the
energy functional
\begin{equation}\label{eq:energy}
  E(t_1, \ldots, t_N)
  = -2 \sum_{i<j} \log|t_i - t_j|
    + \frac{1}{2}\sum_{i=1}^N t_i^2.
\end{equation}
This is the \emph{Dyson logarithmic Coulomb gas}: $N$ unit charges on the
real line with logarithmic repulsion, confined by a quadratic external
potential.

\subsection{Hamiltonian decomposition}

Promoting the positions $t_i$ to canonical coordinates with conjugate
momenta $p_i$ (so $\{t_i, p_j\} = \delta_{ij}$), the energy
\eqref{eq:energy} decomposes into $\binom{N}{2}$ pairwise Hamiltonians:
\begin{equation}\label{eq:H-pair}
  H_{ij} = \frac{p_i^2}{2} + \frac{p_j^2}{2}
           - 2\log|t_i - t_j|
           + \frac{1}{2(N-1)}\bigl(t_i^2 + t_j^2\bigr),
\end{equation}
where the harmonic confinement $\frac{1}{2}\sum t_i^2$ has been
distributed equally among the $N-1$ pairs involving each particle. This
is precisely the structure expected by the \texttt{NBodyAlgebra} engine
with \texttt{potential="log"} and
\texttt{external\_potential=\{"omega": 1\}}.

%% ====================================================================
\section{Mapping to the Poisson Algebra Computation}
\label{sec:mapping}

The \texttt{NBodyAlgebra} class~\cite{3bodyrepo} computes the Lie algebra
$\mathfrak{g}$ generated by the Hamiltonians $\{H_{ij}\}$ under iterated
Poisson brackets:
\begin{align}
  \mathfrak{g}_0 &= \operatorname{span}\{H_{ij}\}, \\
  \mathfrak{g}_k &= \mathfrak{g}_{k-1}
    + \operatorname{span}\bigl\{\{f, H_{ij}\} : f \in \mathfrak{g}_{k-1},
      \, (i,j) \in \text{pairs}\bigr\},
\end{align}
and the dimension sequence $L_k = \dim \mathfrak{g}_k$ is detected via
SVD rank analysis at random phase-space points.

For the GUE log-gas at $N = 3$, $d = 1$, the parameters are:

\begin{center}
\begin{tabular}{ll}
  \hline
  \texttt{n\_bodies} & 3 \\
  \texttt{d\_spatial} & 1 \\
  \texttt{potential} & \texttt{"log"} \\
  \texttt{external\_potential} & \texttt{\{"omega": 1\}} \\
  \texttt{masses} & equal (default) \\
  \texttt{charges} & none (default: gravitational sign) \\
  \hline
\end{tabular}
\end{center}

\subsection{Auxiliary variable trick}

The engine introduces auxiliary variables $u_{ij} = 1/r_{ij}$ where
$r_{ij} = |t_i - t_j|$ (in 1D, $r_{ij} = |t_i - t_j|$). For the
logarithmic potential, $V_{ij} = -\log(r_{ij}) = \log(u_{ij})$, and the
chain rule yields polynomial derivatives in the extended phase space
$(t_i, p_i, u_{ij})$.

%% ====================================================================
\section{Universality and the Dimension Sequence}
\label{sec:universality}

The central empirical result of the parent project is:

\begin{conjecture}[Universality]
  For $N = 3$ bodies interacting via any singular pairwise potential
  $V(r) \sim r^{-\alpha}$ (including $\alpha \to 0$, the logarithmic
  case), in any spatial dimension $d \in \{1, 2, 3\}$, and for any
  positive mass ratios, the dimension sequence is
  \[
    [L_0, L_1, L_2, L_3] = [3, 6, 17, 116].
  \]
\end{conjecture}

This has been verified numerically for:
\begin{itemize}
  \item Potentials: $1/r$, $1/r^2$ (Calogero--Moser pairwise part),
    $1/r^3$, $\log(r)$
  \item Dimensions: $d = 1, 2, 3$
  \item Mass ratios: continuous sweep over $m_2/m_1 \in (0, 5]$
  \item Charge configurations: all-attractive, all-repulsive, mixed
\end{itemize}

The \emph{prediction} for the GUE composite (log + harmonic) is therefore
$[3, 6, 17, 116]$. Confirming this would directly link the algebraic
structure of zeta zero correlations to this universality class.

\subsection{What the dimensions count}

At bracket level $k$, the algebra $\mathfrak{g}_k$ contains all
$(k+2)$-point correlation observables that can be built from pairwise
interactions. For zeta zeros:

\begin{itemize}
  \item \textbf{Level 0} (dim 3): The three pairwise log-interactions.
    These encode the two-point correlation \eqref{eq:pair-correlation}.

  \item \textbf{Level 1} (dim 6): Three new generators appear from
    single brackets $\{H_{ij}, H_{jk}\}$. These are three-point
    correlations of zeros.

  \item \textbf{Level 2} (dim 17): Eleven new generators encode
    four-point correlations. These govern, for instance, the variance of
    the number of zeros in short intervals.

  \item \textbf{Level 3} (dim 116): Ninety-nine new generators encode
    five-point correlations and beyond.
\end{itemize}

\subsection{Connection to primes}

The explicit formula connects zeta zeros to the prime-counting function:
\begin{equation}\label{eq:explicit}
  \psi(x) = x - \sum_\rho \frac{x^\rho}{\rho}
           - \log(2\pi) - \frac{1}{2}\log(1 - x^{-2}),
\end{equation}
where $\psi(x) = \sum_{p^k \le x} \log p$ is the Chebyshev function and
the sum runs over nontrivial zeros $\rho$. The $k$-point correlation
functions of the zeros (encoded in $\mathfrak{g}_k$) determine the
statistical behavior of $\psi(x)$ in short intervals, and hence the
fine-grain distribution of primes.

%% ====================================================================
\section{Computational Experiment}
\label{sec:experiment}

We compute four configurations to test universality for the GUE
Hamiltonian:

\begin{enumerate}
  \item[(a)] \textbf{Pure log-gas}: $V_{ij} = -\log|t_i - t_j|$,
    no confinement. Reference: already known to give $[3, 6, 17, 116]$.

  \item[(b)] \textbf{GUE composite}: $V_{ij} = -\log|t_i - t_j|$ plus
    harmonic confinement $\frac{\omega^2}{2}t_i^2$ with $\omega = 1$.
    This is the novel computation.

  \item[(c)] \textbf{Penning trap}: $V_{ij} = 1/|t_i - t_j|$ plus
    harmonic confinement. Known to give $[3, 6, 17, 116]$ in $d = 2$;
    we verify in $d = 1$.

  \item[(d)] \textbf{Harmonic only}: Pure quadratic potential. Known to
    close at dimension 15 (integrable system).
\end{enumerate}

If (b) yields $[3, 6, 17, 116]$, this confirms that the Poisson algebra
of the exact GUE Hamiltonian governing zeta zero statistics belongs to
the same universality class as Newtonian gravity.

%% ====================================================================
\section{Further Directions}
\label{sec:further}

\subsection{$N > 3$ GUE zeros}

For $N$ zeros, the GUE log-gas has $\binom{N}{2}$ pairwise Hamiltonians.
The parent project has established the formula
\begin{equation}\label{eq:L2}
  L_2(N) = \frac{N(4N^2 - 9N + 3)}{2}, \quad N \ge 4,
\end{equation}
with a boundary correction of $-1$ at $N = 3$.  Running $N = 4, 5$
would test whether this formula holds for the GUE composite.

The graph-theoretic decomposition conjectures that new generators at
level $k$ scale as $c_k \binom{N}{k+1}$, with $c_0 = 1$,
$c_2 = 12$, and a tentative $c_3 = 1198$ from $N = 4$ data. GUE
computations at larger $N$ would provide independent verification.

\subsection{Level 4 and beyond}

Level 4 computations for $N = 3$ take days rather than minutes and
would require AWS dispatch. If level 3 confirms universality for the
GUE composite, a level-4 computation would be a high-value target:
it would determine $L_4$ for the system whose correlations govern
prime distribution statistics.

\subsection{Selberg trace formula}

The Selberg trace formula relates the eigenvalues of the Laplacian on the
modular surface $\mathrm{SL}(2, \Z) \backslash \mathbb{H}$ to a sum over
closed geodesics. The lengths of primitive geodesics are $2\log p$ for
primes $p$ (via the norm map on the ring of integers). The geodesic flow
is a Hamiltonian system on $T^*(\mathrm{SL}(2,\Z)\backslash\mathbb{H})$,
and one could investigate whether the Poisson algebra generated by
interactions between ``nearby'' geodesics carries information about
prime gaps.

This is more speculative: the Hamiltonian is that of free geodesic
motion on a quotient surface, not a pairwise interacting $N$-body
system. However, scattering resonances of the Laplacian are related to
zeta zeros (via the Selberg zeta function), providing a second bridge.

\subsection{Rankin--Selberg $L$-function interactions}

Given Dirichlet characters $\chi_1, \chi_2, \chi_3$ modulo $q$, the
Rankin--Selberg convolution $L(s, \chi_i \times \bar\chi_j)$ provides
a natural ``interaction energy'' between $L$-functions. Define
\[
  H_{ij}(s) = -\log|L(s, \chi_i \times \bar\chi_j)|.
\]
This gives a pairwise Hamiltonian structure on the space of $L$-functions.
While the Poisson bracket must be defined carefully (these are not
canonical phase-space coordinates in the standard sense), the algebraic
structure of nested brackets $\{H_{ij}, H_{jk}\}$ etc.\ could reveal
universal features of $L$-function families.

\subsection{Numerical experiments on zeta zero dynamics}

Beyond the algebraic computation, one could:
\begin{enumerate}
  \item Sample actual zeta zeros $\gamma_1, \ldots, \gamma_N$ and
    compute the log-gas energy \eqref{eq:energy}.
  \item Compare the empirical distribution of energies to GUE predictions.
  \item Study how the Poisson bracket generators, evaluated at actual
    zero positions, correlate with prime-counting statistics.
\end{enumerate}
This would close the loop from algebra $\to$ analysis $\to$ number theory.

%% ====================================================================

\begin{thebibliography}{10}

\bibitem{Montgomery1973}
H.~L. Montgomery, \emph{The pair correlation of zeros of the zeta function},
Proc.\ Symp.\ Pure Math.\ \textbf{24} (1973), 181--193.

\bibitem{Odlyzko1987}
A.~M. Odlyzko, \emph{On the distribution of spacings between zeros of the
zeta function}, Math.\ Comp.\ \textbf{48} (1987), 273--308.

\bibitem{Rudnick1996}
Z.~Rudnick and P.~Sarnak, \emph{Zeros of principal $L$-functions and random
matrix theory}, Duke Math.\ J.\ \textbf{81} (1996), 269--322.

\bibitem{Hejhal1994}
D.~A. Hejhal, \emph{On the triple correlation of zeros of the zeta function},
Int.\ Math.\ Res.\ Not.\ \textbf{7} (1994), 293--302.

\bibitem{Dyson1962}
F.~J. Dyson, \emph{Statistical theory of the energy levels of complex systems},
J.\ Math.\ Phys.\ \textbf{3} (1962), 140--175.

\bibitem{3bodyrepo}
B.~Shepp, \texttt{3body-poisson-algebra}, GitHub repository (2026),
\url{https://github.com/bshepp/3body-poisson-algebra}.

\end{thebibliography}

\end{document}
